\begin{textsample}{POS Dim 2 – human – Score 8.00 – es/es\_ef\_anos\_iniciais\_matematica\_1\_p2.txt}  \label{ex:f2_pos_003}
No que se refere a um processo significativo de construção do conhecimento matemático, é essencial que o estudante seja motivado a questionar, formular, testar e validar as suas próprias hipóteses, verificando a adequação da sua resposta à situação-problema proposta, construindo formas de pensar que o \textbf{levem} a refletir e agir de maneira crítica sobre as questões com as quais se depara no dia a dia.
Por esse motivo, é preciso mostrar que as situações apresentadas em sala de aula possuem alguma relação com processos importantes na \textbf{sociedade}, destacando os campos de aplicações da Matemática e suas especificidades.
Nesse sentido, implica -se que o conhecimento matemático é necessário para todos os estudantes da Educação Básica, seja por sua grande aplicação na \textbf{sociedade} contemporânea, ou pelas suas potencialidades na formação de cidadãos críticos, cientes de suas responsabilidades sociais ( BRASIL, 2017 ).
Reafirmando então, que a Matemática é uma ciência \textbf{humana}, fruto das \textbf{necessidades} e preocupações de diferentes culturas, que contribui para solucionar problemas científicos e tecnológicos, impactando a forma do homem atuar no mundo.

% matched lemmas: humano, levar, necessidade, sociedade
\end{textsample}
